% Chapter 13 -- Claim: Diversification Fails When You Need It Most
% Requires: amsmath, amssymb, booktabs, array, tabularx, longtable, graphicx, tcolorbox (breakable, skins), hyperref or url.
% Defines the callout box only if the preamble has not.
\providecolor{fenavy}{HTML}{1B2A6B}
\providecolor{feshade}{HTML}{F3F5FA}
\makeatletter
\@ifundefined{callout}{%
  \newtcolorbox{callout}{breakable, enhanced,
    colback=white, colframe=gray!70, boxrule=0.5pt, arc=1mm,
    borderline west={2.5pt}{0pt}{gray!70},
    left=8pt, right=8pt, top=4pt, bottom=4pt}%
}{}
\makeatother

\providecolor{fenavy}{HTML}{1B2A6B}
\providecolor{feshade}{HTML}{F3F5FA}
\makeatletter
% Shaded box for optional technical material that general readers may skip.
\@ifundefined{technote}{%
  \newtcolorbox{technote}{breakable, enhanced,
    colback=feshade, colframe=fenavy!60, boxrule=0.5pt, arc=1mm,
    left=8pt, right=8pt, top=4pt, bottom=4pt}%
}{}
\makeatother

\chapter{Claim: Diversification Fails When You Need It Most}
\label{ch:10}
\providecommand{\term}[1]{\emph{#1}}

Investors are told to spread their savings across many kinds of investments so that a loss in one is cushioned by the others. Yet many people who lived through a market crash remember the opposite: on the worst days, everything seemed to fall at once. This chapter asks which memory is closer to the truth.

Diversification\index{diversification} is often called the only free lunch in finance. If you combine assets that do not move in lockstep, the portfolio is less risky than the average of its parts, and you give up no expected return\index{return!expected} to get that benefit. Every balanced fund (one that holds both stocks and bonds), target-date fund (one that shifts from stocks toward bonds as a chosen retirement year approaches) and pension portfolio relies on this arithmetic. The crash of October 1987, when stock markets around the world fell together, raised a doubt that has never gone away: perhaps assets start moving together, and \term{correlations}\index{correlation} rise, exactly when diversification is needed most.

Testing this doubt turned out to be harder than it looks. Because a correlation over all months mixes calm and turbulent periods, Longin and Solnik (2001) studied only the largest monthly moves in international stock markets. They found that markets move more closely together in large falls than in large rises. This is the worst pattern an investor could face, because diversification weakens just when losses arrive. Ang and Chen (2002) found the same pattern among U.S. stock portfolios. Forbes and Rigobon (2002), questioning the evidence for contagion after the crises of the 1990s, then pointed out a trap in measuring it. When markets become more volatile, their measured correlation rises even if the link between them has not changed. Part of the apparent ``contagion'', a crisis spreading from one market to another, is therefore a statistical illusion.

Campbell, Sunderam and Viceira (2017) turned to the other pillar of diversification, government bonds. Bonds protect a stock portfolio when their correlation with stocks is negative, because they then rise as stocks fall. The authors showed that this correlation changed from positive to negative around 2000, after inflation became low and stable. In 2022, when inflation returned, stocks and Treasury bonds fell together.


The answer matters to everyone who allocates money across assets. If correlations really rise in crises, risk models built on data from calm periods understate the losses to expect, and investors need assets that protect them specifically in bad times. If the rise is mostly a statistical illusion, the free lunch is safer than its critics claim. 


Thus, the claim on trial\index{claims-on-trial method} in this chapter is: \textbf{in a market crash, correlations go to one, both among stocks and between stocks and bonds, so diversification breaks down just when investors need it most.}

The lab in Section~\ref{sec:13-lab} repeats the tests of Longin and Solnik (2001) and the correction of Forbes and Rigobon (2002), using free U.S., international and industry returns from the Kenneth R. French Data Library. It also follows the stock--bond correlation of Campbell, Sunderam and Viceira (2017) with free bond yields and inflation data from FRED. The goal is to check whether diversification really fails in crises, and which assets still protected investors when it mattered.

By the end of the chapter you will be able to:

\begin{itemize}
  \item explain why the average correlation limits diversification, and measure its benefit with the diversification ratio\index{volatility}\index{diversification!ratio};
  \item correct crisis correlations for higher volatility with the Forbes--Rigobon\index{correlation!Forbes--Rigobon correction} method, and read tail correlations against a normal benchmark\index{correlation!exceedance}\index{tail dependence};
  \item explain why the stock--bond correlation\index{correlation!stock--bond} changes sign with inflation\index{inflation};
  \item judge a portfolio by what it lost in bad months.
\end{itemize}



\section{Foundations: correlation, diversification and how crisis correlations are measured}
\label{sec:13-foundations}

\subsection{Correlation and diversification}

Consider two assets with weights $w_1$ and $w_2$, volatilities $\sigma_1$ and $\sigma_2$, and correlation $\rho$. The variance\index{variance} of the portfolio is
\begin{equation}
  \sigma_p^2 = w_1^2\sigma_1^2 + w_2^2\sigma_2^2 + 2w_1w_2\rho\,\sigma_1\sigma_2.
  \label{eq:13-var2}
\end{equation}
The last term carries the correlation, so the lower the correlation, the lower the portfolio's risk. Correlation runs from $-1$, when two assets always move in opposite directions, through 0, when their movements are unrelated, to $+1$, when they always move together. 

Now consider $N$ assets that all have the same volatility $\sigma$ and the same correlation $\rho$ with each other. An equally weighted portfolio of them has variance
\begin{equation}
  \sigma_p^2 = \sigma^2\left[\rho + \frac{1 - \rho}{N}\right].
  \label{eq:13-varN}
\end{equation}
As more assets are added, the second term shrinks, and the variance falls toward $\rho\sigma^2$ but never below it. The average correlation therefore sets a floor on how much risk diversification can remove (Figure~\ref{fig:13-floor}). The benefit is often measured by the \term{diversification ratio}\index{diversification!ratio|textbf}: the weighted average volatility of the assets divided by the volatility of the portfolio. A ratio of 1 means no benefit at all, and higher values mean more.

\textbf{Worked example.} Ten industries each have a volatility of 20\%. If their average correlation is 0.6, an equally weighted portfolio has a volatility of $20\% \times \sqrt{0.6 + 0.4/10} = 16.0\%$. If the correlation rises to 0.8, the portfolio's volatility rises to $20\% \times \sqrt{0.8 + 0.2/10} = 18.1\%$. Diversification first cut risk by 4 points, from 20\% to 16\%. After the rise in correlation it cuts risk by less than 2 points, so half of the benefit is gone.

\begin{figure}[htbp]
\centering
\scalebox{0.8}{%
\begin{tikzpicture}[
  % Equally weighted portfolio of N assets, each with 20% volatility
  declare function={pvol(\n,\r) = 20*sqrt(1/\n + (1 - 1/\n)*\r);}
]
\begin{axis}[
  width=0.54\linewidth, height=0.34\linewidth,
  axis lines=left, axis line style={-stealth},
  xmin=1, xmax=50, ymin=14, ymax=20.5,
  xtick={5,10,...,50}, ytick={14,16,18,20},
  xlabel={Number of assets, $N$},
  ylabel={Portfolio volatility (\%)},
  samples at={1,...,50}, clip=false,
]
  % Floors set by the average correlation
  \addplot[densely dotted, domain=1:50, samples=2] {20*sqrt(0.8)};
  \addplot[densely dotted, domain=1:50, samples=2] {20*sqrt(0.6)};
  \addplot[fegray, very thick, dashed] {pvol(x,0.8)};
  \addplot[fenavy, very thick] {pvol(x,0.6)};
  \fill[fegray] (axis cs:10,{pvol(10,0.8)}) circle (2.4pt);
  \fill[fenavy] (axis cs:10,{pvol(10,0.6)}) circle (2.4pt);
  \node[anchor=south west] at (axis cs:10.3,{pvol(10,0.8)+0.05}) {18.1\%};
  \node[anchor=south west] at (axis cs:10.3,{pvol(10,0.6)+0.05}) {16.0\%};
  \node[fegray, anchor=west, align=left] at (axis cs:51,17.9) {$\rho = 0.8$:\\floor 17.9\%};
  \node[fenavy, anchor=west, align=left] at (axis cs:51,15.5) {$\rho = 0.6$:\\floor 15.5\%};
\end{axis}
\end{tikzpicture}%
}
\caption{Adding more assets reduces risk only down to a floor set by how closely they move together. Volatility of an equally weighted portfolio of assets with 20\% volatility each. Adding assets helps only down to a floor set by the average correlation: 15.5\% at a correlation of 0.6 and 17.9\% at 0.8. With ten assets, as in the worked example, a rise from 0.6 to 0.8 lifts volatility from 16.0\% to 18.1\%.}
\label{fig:13-floor}
\end{figure}

\subsection{Why measured correlation rises when volatility rises}

Suppose the return of one market, $y$, is linked to the return of another, $x$, by $y = \beta x + e$, where $e$ is unrelated noise. The link $\beta$ and the size of the noise stay the same, but $x$ becomes more volatile. The correlation between the two markets is
\begin{equation}
  \rho = \frac{\beta\sigma_x}{\sqrt{\beta^2\sigma_x^2 + \sigma_e^2}},
  \label{eq:13-bias}
\end{equation}
which rises as $\sigma_x$ rises. The measured correlation goes up even though the link between the markets has not changed (Figure~\ref{fig:13-bias}). The reason is that the common part of the returns grows while the noise does not, so the common part dominates. Correlations measured in turbulent periods are therefore \term{biased upward} (Boyer, Gibson and Loretan 1999). An everyday picture helps: when a storm hits a whole town, every house gets wet at the same time, but that does not mean the houses have suddenly become more connected to one another.

\begin{figure}[htbp]
\centering
\scalebox{0.8}{%
\begin{tikzpicture}[
  % y = x + e with sigma_e = 20%: measured correlation = sigma_x / sqrt(sigma_x^2 + 20^2)
  declare function={rho(\s) = \s/sqrt(\s*\s + 400);}
]
\begin{axis}[
  width=0.58\linewidth, height=0.34\linewidth,
  axis lines=left, axis line style={-stealth},
  xmin=0, xmax=42, ymin=0, ymax=1,
  xtick={0,5,...,40}, ytick={0,0.2,...,1},
  yticklabel style={/pgf/number format/fixed},
  xlabel={Volatility of market $x$, $\sigma_x$ (\%)},
  ylabel={Measured correlation of $x$ and $y$},
  samples=150, domain=2:42, clip=false,
]
  \addplot[fenavy, very thick] {rho(x)};
  \fill (axis cs:15,0.6) circle (2.4pt);
  \fill (axis cs:30,{rho(30)}) circle (2.4pt);
  \node[anchor=north west] at (axis cs:15.5,0.58) {calm, $\sigma_x = 15\%$: 0.60};
  \node[anchor=north west] at (axis cs:30.5,{rho(30)-0.02}) {turbulent, $\sigma_x = 30\%$: 0.83};
  \node[anchor=west] at (axis cs:17,0.2) {same link throughout: $y = x + e$, $\sigma_e = 20\%$};
\end{axis}
\end{tikzpicture}%
}
\caption{Turbulence alone makes two markets look more closely linked than they are. The volatility bias in measured correlation. The link between the two markets never changes ($y = x + e$, with noise of 20\%), yet the measured correlation rises from 0.60 when $x$ has 15\% volatility to 0.83 when it has 30\%.}
\label{fig:13-bias}
\end{figure}

\subsection{The Forbes--Rigobon correction}

Forbes and Rigobon (2002) proposed a way to remove this bias. Suppose a correlation $\rho_c$ is measured in a crisis, when the variance of market $x$ is $1 + \delta$ times its normal level. The correlation that would hold at normal variance is
\begin{equation}
  \rho^* = \frac{\rho_c}{\sqrt{1 + \delta\,(1 - \rho_c^2)}}.
  \label{eq:13-fr}
\end{equation}
For example, a correlation of 0.8 measured when variance is four times its normal level ($\delta = 3$) corresponds to $0.8/\sqrt{1 + 3 \times 0.36} = 0.55$ in normal times. The correction assumes that the crisis made market $x$ more volatile but did not change the link between the markets. If the link itself strengthens in a crisis, the correction removes part of a real change.

\subsection{Looking only at the tails}

Selecting only extreme observations also changes a correlation. Take two variables that follow a normal distribution\index{distribution!normal}, and keep only the months in which both fall below some threshold. The correlation within those months is lower than in the full sample, and it falls further as the threshold moves deeper into the tail. \term{Exceedance correlations}\index{correlation!exceedance|textbf} (Longin and Solnik 2001; Ang and Chen 2002) are correlations measured only in months when both returns are beyond a threshold, such as their worst 10\%\index{quantile}. They are compared with what a normal distribution with the same overall correlation would produce. The evidence of a crisis effect is the gap between the two, not the level of the correlation itself.

\subsection{Tail dependence}

In plain words, the question is whether the very worst months of two investments tend to arrive together. The shaded box describes how specialists measure this.

\begin{technote}
\textbf{Technical note: measuring tail dependence.} \term{Tail dependence}\index{tail dependence|textbf} asks a simpler question: when one asset has one of its worst months, how likely is the other to have one of its worst months too? If the assets are independent and ``worst'' means the bottom 5\%, the answer is 5\%. For a normal distribution, the answer approaches that independent level as the tail becomes more extreme, so extreme losses rarely coincide. Distributions with fat tails\index{fat tails}, which are modelled with \term{copulas}\index{copula|textbf} (functions that join the distributions of single assets into a joint distribution) such as the Student $t$ or the Clayton, keep the answer higher, so extreme losses tend to arrive together.
\end{technote}

\subsection{The stock--bond correlation}

Government bonds protect a stock portfolio when bad economic news lowers both expected profits and interest rates, because falling rates raise bond prices just as stocks fall. They fail as protection when news about inflation dominates. Rising inflation lowers bond prices directly, and it lowers stock prices too, because investors then \term{discount} future profits at higher rates\index{discount rate}. Over the decades, the stock--bond correlation has therefore switched between positive and negative.


\section{How it works}
\label{sec:13-how}

Three different things can make diversification look weaker in a crisis, and they mean different things for investors.

\begin{itemize}
  \item \textbf{A measurement effect.} When volatility rises, measured correlation rises even if the underlying link is unchanged. This part is an illusion, and correcting for it shrinks the apparent increase.
  \item \textbf{A common shock.} In a crisis, a single force, such as a global panic, a shortage of funding or a fear of recession, comes to drive most returns. Assets that usually respond to separate news now respond to the same news. This is a real change, and it is the one that hurts investors.
  \item \textbf{A change of regime.} Some relationships reverse when the economic environment changes. The main example is the stock--bond correlation, which is negative when news about growth dominates and positive when news about inflation dominates.
\end{itemize}

What matters to an investor is not the correlation itself but what the portfolio earns in bad times. A portfolio can lose less than stocks in a crash even when correlations rise, as long as some of its assets fall less than others. Diversification fails in the sense that matters to an investor only when every holding falls about as much as the riskiest one.

The forces also differ by asset. International stocks are driven largely by the same global stock market, so their correlation with U.S. stocks is high to begin with and may rise further in sell-offs. Industries within one country are linked even more closely. Treasury bonds and commodities respond to different forces, and their relationship with stocks can change sign.


\section{The claim}
\label{sec:13-claim}

\begin{callout}
\textbf{Claim.} Correlations rise in crises, so diversification fails exactly when investors need it most.
\end{callout}

The claim has two parts. The \textbf{statistical part} says that correlations between risky assets are higher in falling and turbulent markets than in calm ones, by more than the volatility bias and a normal distribution would produce. The \textbf{economic part} says that, as a result, diversified portfolios give little protection in the worst months. The first can be true while the second is false, if assets such as Treasury bonds keep their low correlation with stocks while stocks move closer together.

\section{The original evidence}
\label{sec:13-evidence}

\textbf{International stock markets (Longin and Solnik 2001).} Using methods designed for extreme events on five major markets from 1959 to 1996, Longin and Solnik found that the correlation of large negative returns rises in bear markets but not in bull markets, and that only the large negative returns departed from a normal distribution. They concluded that the market's direction, not volatility alone, drives the increase, and warned that selecting months with large moves can create a rise in correlation that is not real.

\textbf{U.S. stocks (Ang and Chen 2002).} U.S. stock portfolios were much more correlated with the market when it fell than when it rose, by 11.6\% more than a normal distribution would produce in falling markets. Small stocks, value stocks\index{value versus growth} and recent losers\index{winners and losers} showed the biggest gap, and models that switch between calm and turbulent regimes captured it better than GARCH\index{GARCH}.

\textbf{Contagion or interdependence (Forbes and Rigobon 2002).} In the 1987 U.S. crash, the 1994 Mexican peso collapse and the 1997 East Asian crisis, several countries showed \term{contagion}\index{contagion|textbf}, correlations that rose during the crisis. After correcting for higher volatility, almost none did. Markets moved together because of strong links that had existed all along, which Forbes and Rigobon called \term{interdependence}\index{interdependence|textbf}.

\textbf{Evidence for investors (Page and Panariello 2018).} From 1970 to 2017, the correlation between U.S. and other developed stock markets was $-17\%$ in the strongest U.S. rallies and $+87\%$ in the worst 1\% of U.S. sell-offs, with a similar pattern for real estate, corporate and \term{high-yield bonds}\index{credit rating}, emerging markets and hedge funds. Treasury bonds usually moved apart from stocks in sell-offs. Because months were selected by U.S. returns, part of the pattern may come from the tail-selection effect of Section~\ref{sec:13-foundations}.

\textbf{The stock--bond correlation (Campbell, Sunderam and Viceira 2017).} Treasury bonds and stocks moved slightly together on average from 1953 to 2009, strongly in the early 1980s and negatively in the 2000s, especially in the downturns of 2000--02 and 2007--09. The authors traced these swings to how inflation and real interest rates move with the economy.

The evidence supports the statistical part for stocks, with caveats about measurement, and contradicts any simple version of the claim for bonds, whose protection depends on the inflation regime.

\section{Lab~\thechapter: Replicating Longin and Solnik (2001) and Forbes and Rigobon (2002)}
\label{sec:13-lab}

\textbf{Claim on trial\index{claims-on-trial method}:} diversification fails when it is needed most. The lab replicates Longin and Solnik (2001), Forbes and Rigobon (2002) and the stock--bond correlation of Campbell, Sunderam and Viceira (2017). It uses U.S. market excess returns\index{return!excess} from 1926, international developed-market returns from 1990 and 10 industry portfolios from the Kenneth R. French Data Library\index{Kenneth French data library}, and the 10-year Treasury yield, consumer prices and the spot oil price from FRED\index{FRED}. Bond returns are built from yields: each month the lab buys a 10-year bond at par and sells it a month later at the new yield.

\begin{itemize}
  \item \textbf{Part A} measures the correlation of U.S. stocks with international stocks, bonds and oil, and among the ten industries, in calm months, in high-volatility months, after the Forbes--Rigobon correction, and in the worst 10\% of U.S. months (Lab Figures~\ref{ch:10}.1 and~\ref{ch:10}.2).
  \item \textbf{Part B} measures tail dependence in the worst 5\%, 10\% and 20\% of months against a normal distribution with the same correlation.
  \item \textbf{Part C} tracks the rolling 36-month stock--bond correlation by decade, by inflation regime and in 2022 (Lab Figure~\ref{ch:10}.3).
  \item \textbf{Part D} compares 100\% stocks, a 60/40 stock--bond portfolio and an equal mix of U.S. stocks, international stocks, bonds and oil in the worst months, in drawdowns\index{drawdown} and in five crises (Lab Figures~\ref{ch:10}.4 and~\ref{ch:10}.5).
\end{itemize}

\label{sec:13-stress}
\textbf{Robustness checks.} Correlations in falling markets, as in Longin and Solnik, and in volatile markets, as in Forbes and Rigobon, answer different questions, and the lab offers both. The Forbes--Rigobon correction assumes a fixed link between markets, so it removes part of any real strengthening; the raw and corrected correlations are best read as upper and lower bounds. The stock--bond correlation depends on whether the sample includes high inflation: bonds look like strong protection from 2000 to 2021, but not in 1965--1985 or 2022. With about 70 high-volatility months, a correlation of 0.7 has a standard error\index{standard error} of about 0.06, so differences of a few hundredths are noise\index{statistical power|textbf}. The lab ends by asking you to reach your own verdict on the scale of Section~\ref{sec:01-method} and compare it with Section~\ref{sec:13-verdict}.

\section{Where the evidence stands}
\label{sec:13-verdict}

The published studies point to a clear pattern: \textbf{correlations among stock markets and among industries rise in falling markets by more than the volatility bias explains, but diversification across asset classes has usually held up, and the protection that bonds give depends on inflation.} The table summarizes the evidence as of November 2026.

{\small
\begin{longtable}{@{}>{\raggedright\arraybackslash}p{0.19\linewidth}>{\raggedright\arraybackslash}p{0.13\linewidth}>{\raggedright\arraybackslash}p{\dimexpr 0.45\linewidth-6\tabcolsep\relax}>{\raggedright\arraybackslash}p{0.23\linewidth}@{}}
\toprule
Situation & Holds? & Why & Evidence \\
\midrule
\endfirsthead
\multicolumn{4}{@{}l}{\emph{(continued)}} \\
\toprule
Situation & Holds? & Why & Evidence \\
\midrule
\endhead
\midrule
\multicolumn{4}{r@{}}{\emph{(continued on the next page)}} \\
\endfoot
\bottomrule
\endlastfoot
Stock correlations in bear markets & \textbf{Yes} & Correlations of large negative returns rise in bear markets but not in bull markets. & Longin and Solnik (2001); Ang and Chen (2002) \\*
\addlinespace[3pt]
After correcting for higher volatility & \textbf{Partly} & Much of the rise disappears, but methods designed to avoid the bias still find a gap between falling and rising markets. & Forbes and Rigobon (2002); Longin and Solnik (2001) \\
\addlinespace[3pt]
Other risky assets in U.S. sell-offs & \textbf{Yes} & Non-U.S. stocks, real estate, corporate bonds, emerging markets and hedge funds moved much more with U.S. stocks in the worst months. & Page and Panariello (2018) \\
\addlinespace[3pt]
Treasury bonds, 2000--2021 & \textbf{No} & The stock--bond correlation was negative, especially in 2000--02 and 2007--09. & Campbell, Sunderam and Viceira (2017) \\
\addlinespace[3pt]
Treasury bonds when inflation dominates & \textbf{Yes} & The correlation was high in the early 1980s, and stocks and bonds fell together in 2022. & Campbell, Sunderam and Viceira (2017); Dimensional Fund Advisors (2023) \\*
\addlinespace[3pt]
Stock--bond portfolios in the worst months & \textbf{Mostly no} & Bonds usually moved apart from stocks in sell-offs, so mixed portfolios lost much less. & Page and Panariello (2018) \\
\end{longtable}
\addtocounter{table}{-1}% uncaptioned longtable: keep table numbers unchanged
}

The claim is true for the diversification investors rely on least, across different stocks, and only conditionally true for the kind they rely on most, between stocks and bonds.

\textbf{What would change the picture.} The claim would gain support if bonds also failed to protect stocks in crises driven by news about growth. Its statistical part would weaken if stock correlations showed no gap between falling and rising markets after the correction.

\label{sec:13-open}
\textbf{Open questions.} Can the stock--bond correlation be forecast from inflation in time to adjust a portfolio, and how closely do private assets, whose smoothed valuations make them look like good diversifiers, really move with public stocks in a crisis?

\section{What practitioners and investors can use}
\label{sec:13-investors}

Plan with crisis correlations, not average ones, and diversify across risks rather than labels: ten stock funds in different regions are mostly one bet on global stocks. Know what your protection protects against, since long-term Treasuries protect against recessions and deflationary panics, not inflation shocks. Judge a portfolio by its losses rather than by whether correlations rose, and rebalance\index{rebalancing|textbf}, because selling what held up and buying what fell helps in the recovery. Ask a manager what the portfolio lost in 2008, March 2020 and 2022, and which holdings would rise if stocks fell 30\% in a recession or in an inflation shock.

\begin{callout}
\textbf{What Failed: bonds as protection in 2022.} For most of the two decades before 2022, Treasury bonds rose when stocks fell, and the 60/40 stock--bond portfolio became the standard choice. Then inflation reached a forty-year high and the Federal Reserve raised rates rapidly. In 2022 the S\&P 500 fell 18.1\% and the U.S. dollar bond market about 13\%, and according to Dimensional Fund Advisors (2023) a 60/40 portfolio lost 11\% over the year. The correlation had been measured in a period of low, stable inflation, when news about growth dominated. The negative stock--bond correlation belonged to one economic regime, not to all of them.
\end{callout}

% Key terms that this chapter uses without defining them, for readers new to finance.
\providecolor{FEMaroon}{RGB}{128,30,45}
\providecolor{FECream}{RGB}{253,248,238}
\providecolor{FEGold}{RGB}{196,150,60}
\makeatletter
\@ifundefined{keyterms}{%
  \newtcolorbox{keyterms}[1][Key terms in this chapter]{breakable, enhanced,
    colback=FECream, colframe=FEGold, boxrule=0.6pt, arc=2pt, leftrule=3pt,
    left=8pt, right=8pt, top=6pt, bottom=6pt, before skip=12pt, after skip=12pt,
    fonttitle=\bfseries, coltitle=white, colbacktitle=FEMaroon, title={#1}}%
}{}
\makeatother

\begin{keyterms}
\small
\begin{description}[leftmargin=0pt, itemsep=2pt, parsep=0pt, font=\normalfont\bfseries\color{FEMaroon}]
  \item[Discount.] To discount future payments is to convert them into their value today using the return investors require, so a higher rate makes them worth less.
  \item[High-yield bonds.] High-yield bonds, also called junk bonds, are bonds that rating agencies judge more likely to default, so they pay higher interest rates.
\end{description}
\end{keyterms}
